% MAIN TEXT, v8 Section 3. Place after the merged Model section.
% Move the existing rho paragraph and analytical phase diagram to the end.
\section{The inward corner}
\label{sec:inward}
\label{sec:threshold}

A corner criterion is useful only if it exposes the quantities that can
and cannot be changed by productivity. Temporarily allow any
differentiable, strictly increasing, concave $F$. At the corner define
\begin{align}
 C_L&=C_0+a\ell(B),&
 \lambda_L&=F'(C_L)a\ell'(B),\label{eq:local-return}\\
 D_j&=F'(C_L)r_j+S_j\qquad(j=n,h,m).
 \label{eq:external-return}
\end{align}

\begin{proposition}[The common corner criterion]
\label{prop:hub}
Under the model assumptions, a necessary condition for the all-inward
corner to be globally optimal is
$D_j\leq p_j\lambda_L$ for every direction. This condition is also
sufficient when $\ell$ is concave. For $F=F_\eta$, it is equivalent to
\begin{equation}
 \boxed{\frac{p_j}{\Gamma(B)}\geq
 \frac{r_j}{\Omega}
 +S_j\frac{(d_0+\Omega)^\eta}{\Omega}}
 \qquad(j=n,h,m).
 \label{eq:hub}
\end{equation}
Productivity enters this expression only through $\Omega$; $\Gamma$
depends on the production function and the budget, not on $a$.
Under concavity, strict inequalities in every direction make every
optimizer all-inward. Weak inequalities establish an inward optimum
but do not by themselves exclude tied external allocations.
\end{proposition}
\begin{proof}
Let $(v_n,v_h,v_m)=(x,y,w)$. Moving a small amount $v_j>0$ into direction
$j$ and reducing $z$ by $p_jv_j$ gives directional derivative
$D_j-p_j\lambda_L$. A positive derivative contradicts optimality.
If $\ell$ is concave, then $C$ is concave and composition with increasing
concave $F$ preserves concavity. Consequently,
\begin{align}
 U(z,x,y,w)-U(B,0,0,0)
 &\leq\lambda_L(z-B)+\sum_jD_jv_j\nonumber\\
 &\leq\sum_j(D_j-p_j\lambda_L)v_j.
 \label{eq:cornerproof}
\end{align}
This proves sufficiency and the strict statement. For $F_\eta$,
multiplication by $(d_0+\Omega)^\eta>0$ gives
$p_ja\ell'(B)\geq r_j+S_j(d_0+\Omega)^\eta$; division by
$\Omega>0$ gives \cref{eq:hub}.
\end{proof}

\begin{corollary}[The original all-inward criterion]
\label{prop:corner}
For $\ell(z)=\ln(1+z)$ and any differentiable, strictly increasing,
concave $F$, the all-inward corner is optimal exactly when
\begin{equation}
 D_j/p_j\leq\lambda_L\qquad(j=n,h,m).
 \label{eq:corner}
\end{equation}
\end{corollary}
\begin{proof}
The local logarithm is concave, so both directions of
\cref{prop:hub} apply. Its derivatives give precisely
\cref{eq:local-return,eq:external-return} with
$\ell'(B)=1/(1+B)$.
\end{proof}

\begin{corollary}[Adding hard support constraints]
\label{cor:feasibility}
If an additional physical feasibility restriction retains the corner,
the criterion $D_j\leq p_j\lambda_L$ remains sufficient under
concavity, as does \cref{eq:hub} for $F_\eta$. Necessity is guaranteed
when each small budget-transfer direction used in the proof remains
admissible. Without that condition, necessity need not hold.
\end{corollary}
\begin{proof}
An optimum over a set is an optimum over any subset containing it.
The necessity argument requires exactly the indicated feasible
perturbations; without them it need not apply.
\end{proof}

\begin{productioncorollary}[The tight logarithmic floor]
\label{cor:production-growth}
Let $\eta=1$ and let $\ell$ be concave. Then
\begin{equation}
 \Gamma(B)\geq B,
 \label{eq:tight-floor}
\end{equation}
with equality if and only if $\ell$ is linear on $[0,B]$.
Every direction with $S_j>0$ requires
$p_j>S_j\Gamma(B)\geq S_jB$ at any finite $a$.
Conversely, for fixed finite $B,d_0,p_j,r_j,S_j$, some concave local
production function and finite $a>0$ make the corner optimal if and
only if $p_j>S_jB$ in every direction with $S_j>0$.
Linear local production provides a witness. Directions with $S_j=0$
retain their computational-return condition but introduce no residual
floor.
\end{productioncorollary}
\noindent\emph{Proof sketch.}
The tangent to a concave production function at $B$ lies above its
value at zero. Equality forces coincidence with the chord on the whole
interval. For the converse use $\ell(z)=z$ and take
$a>\max_j(r_j+d_0S_j)/(p_j-S_jB)$.
The complete proof, including zero slopes, is in \cref{app:floor-proof}.

\begin{corollary}[Logarithmic productivity thresholds]
\label{lem:production}
\label{prop:threshold}
For $\eta=1$ and concave $\ell$, the exact test is
\begin{equation}
 p_j\geq S_j\frac{\ell(B)}{\ell'(B)}
       +\frac{r_j+d_0S_j}{a\ell'(B)}.
 \label{eq:general-production}
\end{equation}
When $p_j>S_j\Gamma(B)$ in every nontrivial direction,
\begin{equation}
 a\geq a_{\rm crit}^{\ell}(B):=
 \max_{j:r_j+S_j>0}
 \frac{r_j+d_0S_j}{p_j\ell'(B)-S_j\ell(B)}
 \label{eq:acrit-general}
\end{equation}
is necessary and sufficient. The maximum over an empty set is zero.
If a nontrivial denominator is non-positive, no finite $a$ suffices.
For $\ell(z)=\ln(1+z)$, set $L=\ln(1+B)$ and recover
\begin{align}
 H_j&=(1+B)LS_j,& A_j&=(1+B)(r_j+d_0S_j),\label{eq:HA}\\
 p_j&\geq H_j+A_j/a,\label{eq:exact}\\
 a_{\rm crit}&=\max_{j:A_j>0}\frac{A_j}{p_j-H_j}.
 \label{eq:acrit}
\end{align}
Thus the general criterion contains both the original production test
and its logarithmic threshold.
\end{corollary}
\begin{proof}
Expand the $\eta=1$ term in \cref{eq:hub} and use
$\Omega=a\ell(B)$. The numerator in \cref{eq:acrit-general} is
positive exactly in nontrivial directions, so rearrangement is valid
exactly for positive denominators. If it is zero, then $r_j=S_j=0$.
Substitution of $\ell'(B)=1/(1+B)$ gives the last three equations.
\end{proof}

The supporting-hyperplane argument is standard
\citep[chapters 3--5]{BoydVandenberghe2004}. The substantive restriction
is the sharp residual floor: no concave redesign of local production
can reduce it below $S_jB$. Improvements in the level of computation
cannot remove this obstruction at logarithmic value.

% MOVE HERE, unchanged: the paragraph beginning
% "A dimensionless robustness diagnostic is" (eq:rho), followed by
% "Analytical phase boundary" (eq:phase-coordinates) and fig:phase.

\subsection{What the floor compares}
The floor is a relative-value statement. An additional unit of a project
can continue to provide independent information or geographical insurance
while the value of more local computation declines. Raising productivity
increases the level of local computation and its physical marginal product
at once. Under logarithmic value these effects partly cancel: the limiting
local marginal utility is $\ell'(B)/\ell(B)=1/\Gamma(B)$, not infinity.
Consequently a fixed external residual slope can remain competitive even
when remote computational production becomes negligible in the comparison.
This limit follows directly by dividing the logarithmic local marginal
utility by $a$ in numerator and denominator.

The distinction also explains the role of existing computation. Increasing
$C_0$ at fixed new-investment technology does not change the asymptotic
floor, but increases the finite-productivity term $d_0S_j$ in
\cref{eq:general-production}. A civilization with more existing useful
service need not place more marginal value on further service. When all
residual slopes are zero, by contrast, $C_0$ and the choice of increasing
$F$ do not affect the computational allocation: maximizing utility is
then equivalent to maximizing $C$. The pure-computation formula in
\cref{eq:budgetcrit} isolates that case. This is why an argument based
only on the attractiveness of virtual experience omits an independent
resource-allocation comparison.

Neither the floor nor the tipping budget measures a physical maximum
number of worlds. They specify when an additional external option beats
a marginal local use under the chosen objective. A system can be physically
capable of supporting its present workload while economically preferring
more territory. It can also prefer inward investment that fails an
independent support requirement. In an application the service ratios and
the allocation inequalities must therefore be evaluated on the same
resource history, rather than selecting a favourable budget for one and a
favourable efficiency for the other. The joint archived example demonstrates
compatibility on a finite interval; the persistence results determine what
additional conditions are needed beyond it.
